
\section{伪全纯截面}\label{sec:8-3}

我们的下一个目标是考察\BookRecordIndex{T0200}局部哈密顿纤维化的伪全纯截面\BookRecordIndex{T0246}模空间。以下始终设
\[
  \pi:\widetilde M\longrightarrow\Sigma
\]
是闭黎曼曲面 $(\Sigma,\mathrm{dvol}_\Sigma,j_\Sigma)$ 上的\BookRecordIndex{T0200}局部哈密顿纤维化，并配有
定义~\ref{def:8-2-1} 意义下的\BookRecordIndex{T0224}联络形式
$\widetilde\omega\in\Omega^2(\widetilde M)$。我们固定定义~\ref{def:8-2-7} 中的\BookRecordIndex{T0306}竖直
殆复结构 $J=\{J_z\}$，并证明：对保持各纤维上\BookRecordIndex{T0363}辛形式的通有\BookRecordIndex{T0402}正合扰动
$\widetilde\omega$，$\widetilde M$ 的伪全纯截面\BookRecordIndex{T0246}模空间是光滑的。多数证明都与
第~\ref{sec:6-7} 节十分接近，我们只作概述。虽然结果是在一般\BookRecordIndex{T0200}局部哈密顿纤维化的框架
中表述的，读者也可以专门考虑第~\ref{sec:8-1} 节讨论的平凡丛
$\pi:\widetilde M=\Sigma\times M\to\Sigma$。在这种情形，初始\BookRecordIndex{T0224}联络形式应取为
$\omega$ 在显然投影 $\operatorname{pr}:\widetilde M\to M$ 下的拉回
\[
  \widetilde\omega_0:=\operatorname{pr}^*\omega.
\]
此外，\BookRecordIndex{T0306}竖直殆复结构空间
$\mathcal J^{\mathrm{Vert}}(\widetilde M,\pi,\widetilde\omega)$ 可与族
\[
  C^\infty\bigl(\Sigma,\mathcal J_\tau(M,\omega)\bigr)
  \ni J:=\{J_z\}_{z\in\Sigma}
\]
的空间等同。每个这样的族都在 $(\widetilde M,\widetilde\omega_0)$ 上给出\BookRecordIndex{T0088}殆复结构
$\widetilde J$，它在 $T_{z,x}(\Sigma\times M)$ 上等于 $(j_\Sigma,J_z)$。

\paragraph{全纯截面的\BookRecordIndex{T0246}模空间。}
固定\BookRecordIndex{T0306}竖直殆复结构
$J\in\mathcal J^{\mathrm{Vert}}(\widetilde M,\pi,\widetilde\omega)$，并令
$\widetilde J\in\mathcal J(\widetilde M,\pi,\widetilde\omega)$ 为 $\widetilde M$ 上相应的
\BookRecordIndex{T0088}殆复结构。给定同调类 $\widetilde A\in H_2(\widetilde M;\mathbb Z)$，考察\BookRecordIndex{T0246}模空间
\[
  \mathcal M(\widetilde A;\widetilde J)
  :=\bigl\{\widetilde u:\Sigma\to\widetilde M\mid
  \bar\partial_{\widetilde J}(\widetilde u)=0,\
  \pi\circ\widetilde u=\operatorname{id}_\Sigma,\
  [\widetilde u]=\widetilde A\bigr\},
\]
它由表示类 $\widetilde A$ 的 $\widetilde J$-全纯截面
$\widetilde u:\Sigma\to\widetilde M$ 组成。注意，只有当
$\pi_*\widetilde A=[\Sigma]$ 时，这个空间才可能非空。若此条件成立，而
$\widetilde u:\Sigma\to\widetilde M$ 是表示类 $\widetilde A$ 的任意
$\widetilde J$-全纯曲线，则
$\pi\circ\widetilde u:\Sigma\to\Sigma$ 是全纯微分同胚。因此
$\widetilde u\circ(\pi\circ\widetilde u)^{-1}$ 是 $\widetilde J$-全纯截面。由此，\BookRecordIndex{T0246}模空间
$\mathcal M(\widetilde A;\widetilde J)$ 可与如下商空间等同：$\widetilde M$ 中所有表示类
$\widetilde A$ 的 $\widetilde J$-全纯曲线所成空间，对 $\Sigma$ 的全纯自同构群取商。
特别地，若 $\Sigma$ 唯一的这种自同构是恒等映射，这里的记号便与
第~\ref{sec:3-1} 节一致。

以下始终固定\BookRecordIndex{T0306}竖直殆复结构 $J$ 和竖直\BookRecordIndex{T0363}辛形式 $\omega_z$，并考察
\[
  \widetilde\omega_H=\widetilde\omega-d\widetilde\sigma_H
\]
这类\BookRecordIndex{T0224}联络形式的\BookRecordIndex{T0402}正合扰动，如式~\eqref{eq:8-1-2} 与注~\ref{rem:8-2-10} 所述。这里
$H$ 是 $\Sigma$ 上的一形式，其值落在由各纤维上均值为零的光滑函数组成的向量丛中。
用 $\mathcal H$ 表示所有这类\BookRecordIndex{T0157}哈密顿扰动所成的空间。若
$\widetilde M=\Sigma\times M$，则可取
$\widetilde\omega=\widetilde\omega_0$，使每个 $z\in\Sigma$ 都有
$\omega_z=\omega$，并使用第~\ref{sec:8-1} 节的公式。

给定 $H\in\mathcal H$，令
$\widetilde J_H\in\mathcal J(\widetilde M,\pi,\widetilde\omega_H)$ 表示由 $J$ 诱导的
\BookRecordIndex{T0088}殆复结构。由习题~\ref{ex:8-1-5}，$\widetilde M$ 的 $\widetilde J_H$-全纯截面在局部是
扰动\BookRecordIndex{T0209}柯西--黎曼方程之解的图像。此外，注~\ref{rem:8-2-10} 表明，截面
$\widetilde u:\Sigma\to\widetilde M$ 是 $\widetilde J_H$-全纯曲线，当且仅当它满足
方程\sourcefootnote{原书式~(8.3.1) 将扰动项印作 $X_H(u)$，但本节没有定义未加波浪号的
$u$，且紧接着把该一形式定义为 $X_H(\widetilde u)$；据此补正自变量。}
\begin{equation}\label{eq:8-3-1}
  \bar\partial_{\widetilde J}(\widetilde u)
  +X_H(\widetilde u)^{0,1}=0,
\end{equation}
其中
\[
  X_H(\widetilde u)
  \in\Omega^1\bigl(\Sigma,\widetilde u^*T^{\mathrm{Vert}}\widetilde M\bigr)
\]
是取值于\BookRecordIndex{T0308}竖直切丛
$T^{\mathrm{Vert}}\widetilde M\subset T\widetilde M$ 的拉回丛的一形式，它由
\[
  T_z\Sigma\longrightarrow T^{\mathrm{Vert}}_{\widetilde u(z)}\widetilde M,
  \qquad \zeta\longmapsto X_{H_\zeta}(\widetilde u(z))
\]
给出。记
\[
  \mathcal M(\widetilde A;J,H)
  :=\mathcal M(\widetilde A;\widetilde J_H)
\]
为表示类 $\widetilde A$ 的 $\widetilde J_H$-全纯截面\BookRecordIndex{T0246}模空间。我们要证明：给定
$J\in\mathcal J^{\mathrm{Vert}}(\widetilde M,\pi,\widetilde\omega)$，对通有\BookRecordIndex{T0157}哈密顿扰动，
这个\BookRecordIndex{T0246}模空间是光滑有限维流形。相应的正则扰动之集记为
\begin{equation}\label{eq:8-3-2}
  \mathcal H_{\mathrm{reg}}(\widetilde A,J)
  :=\bigl\{H\in\mathcal H\mid
  \widetilde u\in\mathcal M(\widetilde A;J,H)
  \Longrightarrow D_{J,H,\widetilde u}\text{ 满射}\bigr\}.
\end{equation}
这里用
\[
  D_{\widetilde u}=D_{J,H,\widetilde u}:
  \Omega^0\bigl(\Sigma,\widetilde u^*T^{\mathrm{Vert}}\widetilde M\bigr)
  \longrightarrow
  \Omega^{0,1}\bigl(\Sigma,\widetilde u^*T^{\mathrm{Vert}}\widetilde M\bigr)
\]
表示与\BookRecordIndex{T0209}柯西--黎曼方程~\eqref{eq:8-3-1} 相联系的\BookRecordIndex{T0346}线性化算子；它的精确定义将在下文给出。
下面的定理是本节的主要结果，其证明将在本节末尾概述。我们用
$c_1^{\mathrm{Vert}}$ 表示\BookRecordIndex{T0308}竖直切丛的第一 Chern 类。

\begin{theorem}\label{thm:8-3-1}
固定同调类 $\widetilde A\in H_2(\widetilde M;\mathbb Z)$ 和\BookRecordIndex{T0306}竖直殆复结构
$J\in\mathcal J^{\mathrm{Vert}}(\widetilde M,\pi,\widetilde\omega)$。
\begin{enumerate}[label=\textup{(\roman*)},leftmargin=3em]
\item 若 $H\in\mathcal H_{\mathrm{reg}}(\widetilde A,J)$，则
$\mathcal M(\widetilde A;J,H)$ 是光滑流形，维数为
\[
  \dim\mathcal M(\widetilde A;J,H)
  =n(2-2g)+2c_1^{\mathrm{Vert}}(\widetilde A).
\]
它带有自然定向。
\item 集合 $\mathcal H_{\mathrm{reg}}(\widetilde A,J)$ 是 $\mathcal H$ 中的\BookRecordIndex{T0300}剩余集。
\end{enumerate}
\end{theorem}


证明几乎与定理~\ref{thm:3-1-6} 的证明相同，只有两点不同。第一，如
第~\ref{sec:6-7} 节一样，我们只考察截面，所以截面空间 $\mathcal B$ 在
$\widetilde u$ 处的切空间 $T_{\widetilde u}\mathcal B$ 由沿 $\widetilde u$ 的竖直向量场
组成。第二，这里考察的是\BookRecordIndex{T0088}殆复结构的受限集合 $\mathcal J_H$，而不是所有被
$\widetilde\omega$ \BookRecordIndex{T0372}驯服（或与 $\widetilde\omega$ \BookRecordIndex{T0348}相容）的\BookRecordIndex{T0088}殆复结构。不过，
$\mathcal J_H$ 是流形，并且由注~\ref{rem:8-2-10} 中的显式公式容易检验：在每一点
$(z,x)\in M_z\subset\widetilde M$，由扰动变化给出的切向量
$(X_h(\widetilde u))^{0,1}$ 张成竖直切空间 $T_xM_z$。这对应于
引理~\ref{lem:3-2-2}，也正是证明相应\BookRecordIndex{T0268}普适模空间为流形所需的事实；参见
注~\ref{rem:3-2-3} 和下文定理~\ref{thm:8-3-1} 的证明。

\begin{remark}\label{rem:8-3-2}
若光滑截面 $\widetilde u:\Sigma\to\widetilde M$ 对每个
$\zeta\in T_z\Sigma$ 和每个 $\widetilde v\in T_{\widetilde u(z)}\widetilde M$ 都满足
\[
  d\pi(\widetilde u(z))\widetilde v=0
  \quad\Longrightarrow\quad
  \widetilde\omega_H\bigl(d\widetilde u(z)\zeta,\widetilde v\bigr)=0,
\]
就称它相对于 $\widetilde\omega_H$ 是\emph{水平的}。这样的截面对每个\BookRecordIndex{T0306}竖直殆复结构
$J$ 都是 $\widetilde J_H$-全纯的。在许多有趣例子中，\BookRecordIndex{T0315}水平截面是唯一能显式写出的
$\widetilde M$ 的 $\widetilde J_H$-全纯截面，有时可用来计算不变量；例如参见 Seidel
\srcite{362, 363}、Polterovich \srcite{325, 326} 和 Akveld--Salamon \srcite{11} 的工作。

也可以在固定 $H$（从而固定\BookRecordIndex{T0314}水平分布）的条件下，对通有\BookRecordIndex{T0306}竖直殆复结构 $J$ 证明类似于
定理~\ref{thm:8-3-1} 的\BookRecordIndex{T0166}横截性结果；不过，这个论证只适用于非\BookRecordIndex{T0315}水平截面。撇开证明的技术
细节，不难看出为何扰动 $J$ 无法使\BookRecordIndex{T0315}水平截面达到横截。若
$\widetilde u:\Sigma\to\widetilde M$ 是 $\widetilde\omega_H$ 的\BookRecordIndex{T0315}水平截面，而算子
$D_{H,\widetilde u}$ 的指标为负，则 $\widetilde u$ 对每个 $J$ 都是
$\widetilde J_H$-全纯的，但 $D_{H,\widetilde u}$ 绝不可能满射。
\end{remark}

\paragraph{配边。}
一如既往，我们需要证明\BookRecordIndex{T0246}模空间的配边类与选择无关。下面的讨论是标准的；为完整起见仍将
其列出。

设
$H_0\in\mathcal H_{\mathrm{reg}}(\widetilde A,J_0)$、
$H_1\in\mathcal H_{\mathrm{reg}}(\widetilde A,J_1)$。任取从 $J_0$ 到 $J_1$ 的\BookRecordIndex{T0306}竖直殆复
结构光滑同伦
\[
  [0,1]\longrightarrow
  \mathcal J^{\mathrm{Vert}}(\widetilde M,\pi,\widetilde\omega),
  \qquad \lambda\longmapsto J_\lambda.
\]
对从 $H_0$ 到 $H_1$ 的光滑同伦
$[0,1]\to\mathcal H:\lambda\mapsto H_\lambda$，定义
\[
  \mathcal W(\widetilde A;\{J_\lambda,H_\lambda\}_\lambda)
  :=\bigl\{(\lambda,\widetilde u)\mid 0\le\lambda\le1,\
  \widetilde u\in\mathcal M(\widetilde A;J_\lambda,H_\lambda)\bigr\}.
\]
记 $\mathcal H(H_0,H_1)$ 为连接 $H_0$ 与 $H_1$ 的所有\BookRecordIndex{T0157}哈密顿扰动光滑同伦之空间。若
$\{H_\lambda\}_\lambda\in\mathcal H(H_0,H_1)$ 满足端点正则，并且对每个
$(\lambda,\widetilde u)\in
\mathcal W(\widetilde A;\{J_\lambda,H_\lambda\}_\lambda)$ 都有
\[
\begin{split}
  \Omega^{0,1}\bigl(\Sigma,\widetilde u^*T^{\mathrm{Vert}}\widetilde M\bigr)
  ={}&\operatorname{im}D_{J_\lambda,H_\lambda,\widetilde u}\\
  &+\mathbb R\left(
  \frac12(\partial_\lambda J_\lambda)
    d_{H_\lambda}(\widetilde u)j_\Sigma
  +(X_{\partial_\lambda H_\lambda}(\widetilde u))^{0,1}\right),
\end{split}
\]
就称它\emph{（关于 $\widetilde A$ 与 $\{J_\lambda\}_\lambda$）正则}。这里
\[
  d_H(\widetilde u):=\Pi_H(\widetilde u)d\widetilde u:
  T\Sigma\longrightarrow\widetilde u^*T^{\mathrm{Vert}}\widetilde M
\]
是 $d\widetilde u$ 与投影 $\Pi_H$ 的复合；$\Pi_H$ 沿由
$\widetilde\omega-d\widetilde\sigma_H$ 确定的\BookRecordIndex{T0314}水平分布投影到\BookRecordIndex{T0308}竖直切丛。\BookRecordIndex{T0407}正则同伦所成的空间
记为
$\mathcal H_{\mathrm{reg}}(\widetilde A,\{J_\lambda\}_\lambda;H_0,H_1)$。下面的结果与
定理~\ref{thm:8-3-1} 相仿，证明从略。

\begin{theorem}\label{thm:8-3-3}
设 $\widetilde A\in H_2(\widetilde M;\mathbb Z)$，
$H_0\in\mathcal H_{\mathrm{reg}}(\widetilde A,J_0)$，
$H_1\in\mathcal H_{\mathrm{reg}}(\widetilde A,J_1)$，并设
\[
  [0,1]\longrightarrow
  \mathcal J^{\mathrm{Vert}}(\widetilde M,\pi,\widetilde\omega),
  \qquad \lambda\longmapsto J_\lambda
\]
是从 $J_0$ 到 $J_1$ 的\BookRecordIndex{T0306}竖直殆复结构光滑同伦。
\begin{enumerate}[label=\textup{(\roman*)},leftmargin=3em]
\item 若
$\{H_\lambda\}_\lambda\in
\mathcal H_{\mathrm{reg}}(\widetilde A,\{J_\lambda\}_\lambda;H_0,H_1)$，则
$\mathcal W(\widetilde A;\{J_\lambda,H_\lambda\}_\lambda)$ 是带边光滑定向流形，且
\[
  \partial\mathcal W(\widetilde A;\{J_\lambda,H_\lambda\}_\lambda)
  =\mathcal M(\widetilde A;J_1,H_1)
   -\mathcal M(\widetilde A;J_0,H_0).
\]
负号表示反向定向。
\item 集合
$\mathcal H_{\mathrm{reg}}(\widetilde A,\{J_\lambda\}_\lambda;H_0,H_1)$ 是
$\mathcal H(H_0,H_1)$ 中的\BookRecordIndex{T0300}剩余集。
\end{enumerate}
\end{theorem}

为准备定理~\ref{thm:8-3-1} 的证明，现在讨论\BookRecordIndex{T0200}局部哈密顿纤维化伪全纯截面的解析设置。
我们先在\BookRecordIndex{T0308}竖直切丛
$T^{\mathrm{Vert}}\widetilde M\to\widetilde M$ 上定义一个自然联络，再用它导出线性化
\[
  D_{J,H,\widetilde u}:
  \Omega^0\bigl(\Sigma,\widetilde u^*T^{\mathrm{Vert}}\widetilde M\bigr)
  \longrightarrow
  \Omega^{0,1}\bigl(\Sigma,\widetilde u^*T^{\mathrm{Vert}}\widetilde M\bigr)
\]
的显式公式。正如第~\ref{sec:3-1} 节开头所指出的，证明定理~\ref{thm:8-3-1} 实际上
不需要这样的公式：若 $\widetilde u$ 是 $\widetilde J_H$-全纯的，则
$D_{J,H,\widetilde u}$ 不依赖于联络，并可理解为\BookRecordIndex{T0310}竖直微分。只需知道它是\BookRecordIndex{T0210}柯西--黎曼算子，
从而是\BookRecordIndex{T0128}弗雷德霍姆算子。承认这一点，便可不借助下文的显式公式来理解
定理~\ref{thm:8-3-1} 的证明。

\paragraph{\BookRecordIndex{T0312}竖直 Levi--Civita 联络。}
向量丛 $T^{\mathrm{Vert}}\widetilde M\to\widetilde M$ 带有一个由\BookRecordIndex{T0306}竖直殆复结构 $J$ 和
$\widetilde\omega$ 的\BookRecordIndex{T0314}水平分布诱导的联络 $\nabla$。它称为\emph{\BookRecordIndex{T0312}竖直 Levi--Civita
联络}，并由下列公理刻画。

\medskip
\noindent\textbf{（竖直）}\quad
$\nabla$ 在每个纤维 $(M_z,\omega_z)$ 上的限制，是度量
\[
  g_z=\frac12\bigl(\omega_z(\,\cdot\,,J_z\,\cdot\,)
    -\omega_z(J_z\,\cdot\,,\,\cdot\,)\bigr)
\]
的 \BookRecordIndex{T0034}Levi--Civita 联络。

\medskip
\noindent\textbf{（水平）}\quad
设 $\gamma:\mathbb R\to\Sigma$ 是光滑路径，并以
$\psi_t:M_{\gamma(0)}\to M_{\gamma(t)}$ 表示关于 $\widetilde\omega$ 的\BookRecordIndex{T0314}水平分布的\BookRecordIndex{T0263}平行
移动。假设 $(\widetilde x,\xi):\mathbb R\to T^{\mathrm{Vert}}\widetilde M$ 是
$\gamma$ 的\BookRecordIndex{T0317}水平提升，即
\[
  \pi(\widetilde x(t))=\gamma(t),\qquad
  \widetilde x(t)=\psi_t(\widetilde x(0)),\qquad
  \xi(t)=d\psi_t(\widetilde x(0))\xi(0).
\]
则沿 $\widetilde x$ 的竖直向量场 $\xi$ 关于 $\nabla$ 平行。

\medskip
\noindent\textbf{（自然性）}\quad
若 $\psi:\widetilde M_0\to\widetilde M_1$ 是逐纤维微分同胚，且
$\psi^*\widetilde\omega_1=\widetilde\omega_0$、$\psi^*J_1=J_0$，则
$\psi^*\nabla_1=\nabla_0$。

\medskip
\noindent\textbf{（哈密顿）}\quad
设 $\widetilde\sigma_H\in\Omega^1(\widetilde M)$ 是水平一形式。则由 $J$ 与
$\widetilde\omega-d\widetilde\sigma_H$ 所决定的\BookRecordIndex{T0312}竖直 Levi--Civita 联络 $\nabla_H$ 满足
\begin{equation}\label{eq:8-3-3}
  \nabla_{H,\widetilde v}\xi(\widetilde x)
  =\nabla_{\widetilde v}\xi(\widetilde x)
   +\nabla_\xi X_{H_\zeta}(\widetilde x),
  \qquad \zeta:=d\pi(\widetilde x)\widetilde v,
\end{equation}
其中 $\xi:\widetilde M\to T^{\mathrm{Vert}}\widetilde M$ 是任意竖直向量场，
$\widetilde v\in T_{\widetilde x}\widetilde M$。

\BookRecordIndex{T0312}竖直 Levi--Civita 联络由前两条公理唯一确定。换言之，$\nabla$ 只是逐纤维
\BookRecordIndex{T0034}Levi--Civita 联络与纤维丛 $T^{\mathrm{Vert}}\widetilde M\to\Sigma$ 中辛联络的组合。
前两条条件显然在逐纤维微分同胚下不变，故与（自然性）公理\BookRecordIndex{T0348}相容。在平凡丛
$\widetilde M=\Sigma\times M$ 上，也可把（哈密顿）公理视为\BookRecordIndex{T0312}竖直 Levi--Civita 联络的
定义；这时还须验证它满足（水平）公理，而这直接来自注~\ref{rem:8-2-10} 中
$\widetilde\omega_H$ 的\BookRecordIndex{T0314}水平分布公式。明确这一点后，可用（哈密顿）公理在
命题~\ref{prop:8-2-2} 的局部平凡化中定义联络，再用（自然性）公理证明它在整体上良定。

\begin{exercise}\label{ex:8-3-4}
设 $\nabla$ 是 $\widetilde\omega$ 的\BookRecordIndex{T0312}竖直 Levi--Civita 联络，
$H\in\mathcal H$ 是\BookRecordIndex{T0157}哈密顿扰动，并设 $\nabla_H$ 是由式~\eqref{eq:8-3-3} 定义在
$T^{\mathrm{Vert}}\widetilde M\to\widetilde M$ 上的联络。证明：$\nabla_H$ 对
$\widetilde\omega_H$ 满足（水平）公理。提示：考虑
$\widetilde M=\Sigma\times M$、$\widetilde\omega=\operatorname{pr}^*\omega$ 的情形。
\end{exercise}

\begin{exercise}\label{ex:8-3-5}
设 $\nabla$ 是由 $\widetilde\omega$ 与 $J$ 诱导的
$T^{\mathrm{Vert}}\widetilde M$ 上的\BookRecordIndex{T0312}竖直 Levi--Civita 联络。设
\[
  \mathbb R^2\longrightarrow\widetilde M,
  \qquad(s,t)\longmapsto\widetilde x(s,t)
\]
是满足 $d\pi(\widetilde x)\partial_t\widetilde x=0$ 的光滑映射。证明
\[
  \nabla_s\partial_t\widetilde x
  =\nabla_s\bigl(\Pi(\widetilde x)\partial_t\widetilde x\bigr),
\]
其中 $\Pi:T\widetilde M\to T^{\mathrm{Vert}}\widetilde M$ 表示沿\BookRecordIndex{T0314}水平分布
$\operatorname{Hor}$ 投影到竖直子空间的投影（见式~\eqref{eq:8-2-5}）。
\end{exercise}

一般而言，\BookRecordIndex{T0312}竖直 Levi--Civita 联络沿自身并非竖直的方向并不保持竖直内积。与
第~\ref{sec:3-1} 节一样，下一步是定义同时保持\BookRecordIndex{T0088}殆复结构与内积的联络
$\widetilde\nabla$。下面的引理说明仍可使用先前的同一个公式。

\begin{lemma}\label{lem:8-3-6}
设 $(\widetilde M,\pi,\widetilde\omega)$ 是黎曼曲面
$(\Sigma,j_\Sigma,\mathrm{dvol}_\Sigma)$ 上的\BookRecordIndex{T0200}局部哈密顿纤维化，且
$J\in\mathcal J^{\mathrm{Vert}}(\widetilde M,\pi,\widetilde\omega)$。令 $\nabla$ 为
$T^{\mathrm{Vert}}\widetilde M$ 上相应的\BookRecordIndex{T0312}竖直 Levi--Civita 联络。则公式
\[
  \widetilde\nabla_{\widetilde v}X
  :=\nabla_{\widetilde v}X-\frac12J(\nabla_{\widetilde v}J)X,
  \qquad X\in\Omega^0(\widetilde M,T^{\mathrm{Vert}}\widetilde M),
\]
在 $T^{\mathrm{Vert}}\widetilde M$ 上定义了一个\BookRecordIndex{T0052}埃尔米特联络。
\end{lemma}

\begin{proof}
只须考虑平凡丛 $\mathbb R\times M$。此时\BookRecordIndex{T0306}竖直殆复结构是光滑映射
\[
  \mathbb R\longrightarrow\mathcal J(M,\omega),
  \qquad t\longmapsto J_t,
\]
且
\[
  \widetilde\omega_H=\omega-dH\wedge dt,
\]
其中 $H:\mathbb R\times M\to\mathbb R$。在方便时记 $H_t(x):=H(t,x)$，并用
$X_t:=X_{H_t}\in\Vect(M,\omega)$ 表示 $H_t$ 的\BookRecordIndex{T0160}哈密顿向量场。考察路径
\[
  \mathbb R\longrightarrow\mathbb R\times M,
  \qquad t\longmapsto(t,x(t))
\]
以及沿它的竖直向量场 $\xi(t)\in T_{x(t)}M$。于是
\[
  \nabla_{H,t}\xi=\nabla_t\xi+\nabla_\xi X_t(x),
\]
其中 $\nabla=\nabla^t$ 表示 $TM$ 上度量 $\omega(\,\cdot\,,J_t\,\cdot\,)$ 的
\BookRecordIndex{T0034}Levi--Civita 联络。直接计算可得
\[
  \nabla_{H,t}J(x)
  =\nabla_{\dot x+X_t(x)}J_t(x)+(\partial_tJ_t)(x)
   -\mathcal L_{X_t}J_t(x),
\]
这里
\[
  (\mathcal L_XJ)\xi
  =(\nabla_XJ)\xi-\nabla_{J\xi}X+J\nabla_\xi X
\]
表示 $J$ 沿 $X$ 方向的 Lie 导数。因此
\begin{align*}
  \widetilde\nabla_{H,t}\xi
  &=\nabla_{H,t}\xi-\frac12J(\nabla_{H,t}J)\xi\\
  &=\nabla_t\xi-\frac12J(\partial_tJ)\xi
    -\frac12J(\nabla_{\dot x+X}J)\xi
    +\nabla_\xi X+\frac12J(\mathcal L_XJ)\xi.
\end{align*}
这里略去了下标 $t$ 和自变量 $x$。公式
\[
  \frac d{dt}\langle\xi_1,\xi_2\rangle
  =\langle\nabla_t\xi_1,\xi_2\rangle
   +\langle\xi_1,\nabla_t\xi_2\rangle
   -\langle\xi_1,J(\partial_tJ)\xi_2\rangle
\]
表明 $\nabla_t-\frac12J(\partial_tJ)$ 是黎曼联络。此外，
$T_xM$ 的自同态 $J(\nabla_{\dot x+X}J)$ 是\BookRecordIndex{T0350}斜伴随的。公式
\[
  \langle\xi_1,\nabla_{\xi_2}X\rangle
  +\langle\nabla_{\xi_1}X,\xi_2\rangle
  +\langle\xi_1,J(\mathcal L_XJ)\xi_2\rangle=0
\]
（对\BookRecordIndex{T0160}哈密顿向量场成立）说明自同态
\[
  \xi\longmapsto\nabla_\xi X+\frac12J(\mathcal L_XJ)\xi
\]
也是\BookRecordIndex{T0350}斜伴随的。因此 $\widetilde\nabla_H$ 保持黎曼度量；它保持\BookRecordIndex{T0088}殆复结构则直接来自定义。
\end{proof}

\paragraph{\BookRecordIndex{T0346}线性化算子。}
现在可以引入 $\widetilde J$-全纯截面的\BookRecordIndex{T0346}线性化算子。固定无边紧黎曼曲面
$(\Sigma,j_\Sigma,\mathrm{dvol}_\Sigma)$、带\BookRecordIndex{T0224}联络形式 $\widetilde\omega$ 的\BookRecordIndex{T0200}局部哈密顿
纤维化 $\pi:\widetilde M\to\Sigma$，以及\BookRecordIndex{T0306}竖直殆复结构
$J\in\mathcal J^{\mathrm{Vert}}(\widetilde M,\pi,\widetilde\omega)$。对任意光滑截面
$\widetilde u:\Sigma\to\widetilde M$，定义线性算子
\[
  D_{\widetilde u}=D_{J,\widetilde u}:
  \Omega^0\bigl(\Sigma,\widetilde u^*T^{\mathrm{Vert}}\widetilde M\bigr)
  \longrightarrow
  \Omega^{0,1}\bigl(\Sigma,\widetilde u^*T^{\mathrm{Vert}}\widetilde M\bigr)
\]
为
\[
  D_{\widetilde u}\xi
  :=(\nabla\xi)^{0,1}
   -\frac12J(\nabla_\xi J)\partial_{\widetilde J}(\widetilde u),
\]
其中 $\nabla$ 表示\BookRecordIndex{T0312}竖直 Levi--Civita 联络，$\widetilde J$ 表示由
$\widetilde\omega$ 与 $J$ 诱导的 $\widetilde M$ 上\BookRecordIndex{T0088}殆复结构。与第~\ref{sec:3-1} 节一样，
这是实线性\BookRecordIndex{T0210}柯西--黎曼算子。若 $J$ 与 $\widetilde\omega$ 属于 $C^\ell$ 类，而
$\widetilde u$ 属于 $W^{\ell,p}$ 类，其中 $\ell$ 是整数、$p>2$，则对每个
$k\in\{1,\ldots,\ell\}$，$D_{\widetilde u}$ 都良定地把
$\widetilde u^*T^{\mathrm{Vert}}\widetilde M$ 的 $W^{k,p}$ 截面映到 $\Sigma$ 上取值于该
拉回丛的 $W^{k-1,p}$ 类 $(0,1)$-形式。它是 $W^{\ell-1,p}$ 类的实线性\BookRecordIndex{T0210}柯西--黎曼算子，
所以定理 C.2.3 适用。因此，命题~\ref{prop:3-1-11} 推广到算子
\[
  D_{\widetilde u}:
  W^{k,p}\bigl(\Sigma,\widetilde u^*T^{\mathrm{Vert}}\widetilde M\bigr)
  \longrightarrow
  W^{k-1,p}\bigl(\Sigma,
  \Lambda^{0,1}T^*\Sigma\otimes_J
  \widetilde u^*T^{\mathrm{Vert}}\widetilde M\bigr).
\]
特别地，该算子是弗雷德霍姆的，且
\[
  \operatorname{index}D_{\widetilde u}
  =n(2-2g)+2\bigl\langle c_1(T^{\mathrm{Vert}}\widetilde M),[\widetilde u]\bigr\rangle.
\]

\begin{remark}\label{rem:8-3-7}
引入这个算子的几何动机来自与第~\ref{sec:3-1} 节类似的抽象设置。对任意截面
$\widetilde u:\Sigma\to\widetilde M$，无论它是否 $\widetilde J$-全纯，都有映射
\[
  \mathcal F_{\widetilde u}:
  \Omega^0\bigl(\Sigma,\widetilde u^*T^{\mathrm{Vert}}\widetilde M\bigr)
  \longrightarrow
  \Omega^{0,1}\bigl(\Sigma,\widetilde u^*T^{\mathrm{Vert}}\widetilde M\bigr),
\]
定义为
\[
  \mathcal F_{\widetilde u}(\xi)
  :=\Phi_{\widetilde u}(\xi)^{-1}
  \bar\partial_{\widetilde J}\bigl(\exp_{\widetilde u}(\xi)\bigr).
\]
这里，$\exp_{\widetilde x}:T_{\widetilde x}M_z\to M_z$ 表示由 $\omega_z$ 与 $J_z$ 确定的
度量 $g_z$ 所对应的竖直指数映射，而
\[
  \Phi_{\widetilde x}(\xi):T_{\widetilde x}M_z
  \longrightarrow T_{\exp_{\widetilde x}(\xi)}M_z
\]
表示关于 $\widetilde\nabla$ 沿竖直测地线的\BookRecordIndex{T0263}平行移动。与命题~\ref{prop:3-1-1} 一样，
$\mathcal F_{\widetilde u}$ 在零点的微分为
$d\mathcal F_{\widetilde u}(0)=D_{\widetilde u}$。
\end{remark}

\begin{remark}\label{rem:8-3-8}
设 $\widetilde\sigma_H\in\Omega^1(\widetilde M)$ 是注~\ref{rem:8-2-10} 中的水平一形式。
把相应的非线性\BookRecordIndex{T0210}柯西--黎曼算子，即式~\eqref{eq:8-3-1} 的左端，记为
\[
  \bar\partial_{J,H}(\widetilde u)
  :=\bar\partial_{\widetilde J}(\widetilde u)
    +(X_H(\widetilde u))^{0,1}_J.
\]
则与 $J$ 和 $\widetilde\omega-d\widetilde\sigma_H$ 相联系的\BookRecordIndex{T0346}线性化算子
\[
  D_{H,\widetilde u}=D_{J,H,\widetilde u}:
  \Omega^0\bigl(\Sigma,\widetilde u^*T^{\mathrm{Vert}}\widetilde M\bigr)
  \longrightarrow
  \Omega^{0,1}\bigl(\Sigma,\widetilde u^*T^{\mathrm{Vert}}\widetilde M\bigr)
\]
由下式给出：
\begin{align*}
  D_{H,\widetilde u}\xi
  &=(\nabla_H\xi)^{0,1}
    -\frac12J(\nabla_\xi J)\partial_{J,H}(\widetilde u)\\
  &=D_{\widetilde u}\xi
    +\left(\nabla_\xi X_H(\widetilde u)
      -\frac12J(\nabla_\xi J)X_H(\widetilde u)\right)^{0,1}.
\end{align*}
这里 $\nabla_H$ 表示由\BookRecordIndex{T0224}联络形式
$\widetilde\omega-d\widetilde\sigma_H$ 和\BookRecordIndex{T0306}竖直殆复结构 $J$ 所决定的
$T^{\mathrm{Vert}}\widetilde M$ 上的\BookRecordIndex{T0312}竖直 Levi--Civita 联络。
\end{remark}

\begin{proof}[定理~\ref{thm:8-3-1} 的证明]
由定义和\BookRecordIndex{T0380}隐函数定理 A.3.3，对每个
$H\in\mathcal H_{\mathrm{reg}}(\widetilde A,J)$，
$\mathcal M(\widetilde A;J,H)$ 都是具有所述维数的流形。其自然定向与
定理~\ref{thm:3-1-6}\textup{(i)} 的证明相同。

为证明 (ii)，固定整数 $\ell\ge2$，用 $\mathcal H^\ell$ 表示 $C^\ell$ 类\BookRecordIndex{T0157}哈密顿扰动所成
的巴拿赫空间，并考察\BookRecordIndex{T0268}普适模空间
\[
  \mathcal M(\widetilde A;J,\mathcal H^\ell)
  :=\bigl\{(\widetilde u,H)\mid H\in\mathcal H^\ell,\
  \widetilde u\in\mathcal M(\widetilde A;J,H)\bigr\}.
\]
这个空间是\BookRecordIndex{T0053}巴拿赫流形 $\mathcal B^{1,p}\times\mathcal H^\ell$ 上一个
$C^{\ell-1}$ 巴拿赫空间丛的 $C^{\ell-1}$ 截面的零点集；这里
$\mathcal B^{1,p}$ 表示所有表示类 $\widetilde A$ 的 $W^{1,p}$ 截面
$\widetilde u:\Sigma\to\widetilde M$ 所成的空间。为证明这个\BookRecordIndex{T0268}普适模空间本身是
$C^{\ell-1}$ \BookRecordIndex{T0053}巴拿赫流形，我们会遇到\BookRecordIndex{T0346}线性化算子
\[
  W^{1,p}\bigl(\Sigma,\widetilde u^*T^{\mathrm{Vert}}\widetilde M\bigr)
  \times\mathcal H^\ell
  \longrightarrow
  L^p\bigl(\Sigma,\Lambda^{0,1}T^*\Sigma\otimes_J
    \widetilde u^*T^{\mathrm{Vert}}\widetilde M\bigr),
\]
它由
\[
  (\xi,h)\longmapsto D_{H,\widetilde u}\xi+(X_h(\widetilde u))^{0,1}
\]
给出。必须证明：对每个 $H\in\mathcal H^\ell$ 和每个
$\widetilde u\in\mathcal M(\widetilde A;J,H)$，这个算子都满射。

为此，取 $q>1$ 使 $1/p+1/q=1$，并设截面
\[
  \eta\in L^q\bigl(\Sigma,\Lambda^{0,1}T^*\Sigma\otimes_J
  \widetilde u^*T^{\mathrm{Vert}}\widetilde M\bigr)
\]
消去该算子的像。由于 $\widetilde J_H$ 是 $C^{\ell-1}$ 类\BookRecordIndex{T0088}殆复结构，由
命题~\ref{prop:3-1-10}，$\widetilde u$ 属于 $W^{\ell,p}$ 类，从而
$D_{H,\widetilde u}$ 是 $W^{\ell-1,p}$ 类的实线性\BookRecordIndex{T0210}柯西--黎曼算子。因此由定理 C.2.3，
$\eta\in W^{\ell,q}$\sourcefootnote{原书把\BookRecordIndex{T0329}椭圆正则性提升后的空间印作
$W^{\ell,p}$；但 $\eta$ 起初属于 $L^q$，其中 $q$ 是 $p$ 的对偶指数，对伴随算子
$D^*_{H,\widetilde u}$ 应用\BookRecordIndex{T0329}椭圆正则性时积分指数仍为 $q$。这与第~\ref{sec:6-7} 节
同型论证中的指数一致。}，且
\[
  D^*_{H,\widetilde u}\eta=0,
\]
并且对每个 $h\in\mathcal H^\ell$ 都有
\begin{equation}\label{eq:8-3-4}
  \int_\Sigma\bigl\langle\eta,X_h(\widetilde u)\bigr\rangle
  \,\mathrm{dvol}_\Sigma=0.
\end{equation}
我们证明 $\eta$ 必须为零。否则，存在 $z_0\in\Sigma$ 使 $\eta(z_0)\ne0$。于是存在
线性映射
\[
  T_{z_0}\Sigma\longrightarrow C_0^\infty(M_{z_0}),
  \qquad \zeta\longmapsto h_{0,\zeta},
\]
使得
\[
  \bigl\langle\eta(z_0),X_{h_0}(\widetilde u(z_0))\bigr\rangle>0.
\]
任取在 $z=z_0$ 处与 $h_0$ 相同的\BookRecordIndex{T0157}哈密顿扰动 $h\in\mathcal H^\ell$。那么标量函数
$\langle\eta,X_h(\widetilde u)\rangle$ 在 $z_0$ 的某个邻域
$U_0\subset\Sigma$ 中为正。选取支集位于 $U_0$ 且 $\beta(z_0)=1$ 的光滑\BookRecordIndex{T0188}截断函数
$\beta:\Sigma\to[0,1]$。把式~\eqref{eq:8-3-4} 中的 $h$ 换成 $\beta h$ 后，左端积分为
正，矛盾。因此 $\eta$ 必须消失，普适\BookRecordIndex{T0346}线性化算子在每个
$(\widetilde u,H)\in\mathcal M(\widetilde A;J,\mathcal H^\ell)$ 处都满射。

由\BookRecordIndex{T0380}隐函数定理 A.3.3，$\mathcal M(\widetilde A;J,\mathcal H^\ell)$ 是
$C^{\ell-1}$ \BookRecordIndex{T0053}巴拿赫流形。作为\BookRecordIndex{T0213}可分\BookRecordIndex{T0053}巴拿赫流形的子流形，它本身也\BookRecordIndex{T0213}可分，即含有可数稠密
子集。现在把 萨德--斯梅尔定理应用于投影
\[
  \mathcal M(\widetilde A;J,\mathcal H^\ell)\longrightarrow\mathcal H^\ell.
\]
余下论证与定理~\ref{thm:3-1-6} 完全相同，故从略。这就证明了
定理~\ref{thm:8-3-1}。
\end{proof}
